%==============================================================================
%  Evaluacion Corta 2 -- Diseno de Alto Nivel (MP6160)
%  Reporte tecnico: HLS + Prototipo Virtual
%
%  ESTRUCTURA:
%    Parte I  -> Implementacion en HLS      (PENDIENTE: la completa el equipo de HLS)
%    Parte II -> Prototipo Virtual          (COMPLETA)
%    Parte III-> Correspondencia y verificacion cruzada (parcialmente pendiente)
%
%  COMPILAR:   pdflatex reporte_MP6160.tex   (dos veces, para el indice)
%           o  latexmk -pdf reporte_MP6160.tex
%==============================================================================
\documentclass[11pt,letterpaper]{article}

%------------------------- Codificacion e idioma ------------------------------
\usepackage[utf8]{inputenc}
\usepackage[T1]{fontenc}
\usepackage[spanish,es-noquoting]{babel}
\usepackage{lmodern}

%------------------------------ Geometria -------------------------------------
\usepackage[margin=2.5cm]{geometry}
\usepackage{setspace}
\onehalfspacing
% Margen extra de elasticidad: evita desbordes por identificadores largos
% en \texttt{} que no admiten division silabica.
\setlength{\emergencystretch}{3em}

%------------------------------- Graficos -------------------------------------
\usepackage{graphicx}
\usepackage{float}
\usepackage{xcolor}

%-------------------------------- Tablas --------------------------------------
\usepackage{booktabs}
\usepackage{array}
\usepackage{longtable}
\usepackage{multirow}

%-------------------------------- Codigo --------------------------------------
\usepackage{listings}

%------------------------------ Utilidades ------------------------------------
\usepackage{fancyhdr}
\usepackage{enumitem}
\usepackage{amsmath}
\usepackage{amssymb}
\usepackage{caption}
\usepackage[colorlinks=true,linkcolor=azulunc,citecolor=azulunc,%
            urlcolor=azulunc]{hyperref}

%------------------------------- Colores --------------------------------------
\definecolor{azulunc}{HTML}{0F4C81}
\definecolor{verdeok}{HTML}{2E7D32}
\definecolor{rojoerr}{HTML}{B00020}
\definecolor{grisfondo}{HTML}{F2F3F5}
\definecolor{griscom}{HTML}{6A737D}
\definecolor{amarpen}{HTML}{FFF8E1}
\definecolor{bordepen}{HTML}{F9A825}

%--------------------------- Estilo de listings -------------------------------
\lstdefinestyle{base}{
  basicstyle=\ttfamily\footnotesize,
  backgroundcolor=\color{grisfondo},
  commentstyle=\color{griscom}\itshape,
  keywordstyle=\color{azulunc}\bfseries,
  stringstyle=\color{verdeok},
  numberstyle=\tiny\color{griscom},
  breaklines=true,
  breakatwhitespace=false,
  showstringspaces=false,
  tabsize=4,
  frame=leftline,
  framerule=1.2pt,
  rulecolor=\color{azulunc},
  xleftmargin=10pt,
  framexleftmargin=10pt,
  aboveskip=10pt,
  belowskip=10pt,
  captionpos=b,
  literate={á}{{\'a}}1 {é}{{\'e}}1 {í}{{\'i}}1 {ó}{{\'o}}1 {ú}{{\'u}}1
           {Á}{{\'A}}1 {É}{{\'E}}1 {Í}{{\'I}}1 {Ó}{{\'O}}1 {Ú}{{\'U}}1
           {ñ}{{\~n}}1 {Ñ}{{\~N}}1 {ü}{{\"u}}1 {°}{{$^\circ$}}1
}
\lstset{style=base}
\lstdefinestyle{consola}{style=base,language=,rulecolor=\color{verdeok}}
% Estilo para diagramas ASCII anchos (mas de 78 columnas)
\lstdefinestyle{diagrama}{style=base,language=,basicstyle=\ttfamily\scriptsize}
\lstdefinestyle{cstyle}{style=base,language=C}
\lstdefinestyle{cppstyle}{style=base,language=C++}
\lstdefinestyle{pystyle}{style=base,language=Python}
\lstdefinestyle{shstyle}{style=base,language=bash}

%---------------------- Caja para bloques pendientes --------------------------
\newcommand{\pendiente}[1]{%
  \begin{center}
  \fcolorbox{bordepen}{amarpen}{%
    \begin{minipage}{0.93\textwidth}
      \vspace{2pt}
      \textbf{\textcolor{bordepen}{$\blacktriangleright$ PENDIENTE:}} #1
      \vspace{2pt}
    \end{minipage}}
  \end{center}
  \vspace{4pt}
}

%------------------------------ Encabezados -----------------------------------
\pagestyle{fancy}
\fancyhf{}
\fancyhead[L]{\small\textcolor{griscom}{MP6160 -- Evaluación Corta 2}}
\fancyhead[R]{\small\textcolor{griscom}{HLS y Prototipos Virtuales}}
\fancyfoot[C]{\small\thepage}
\renewcommand{\headrulewidth}{0.4pt}

%==============================================================================
\begin{document}
%==============================================================================

%------------------------------- PORTADA --------------------------------------
\begin{titlepage}
\centering
\vspace*{2cm}

{\large\scshape Instituto Tecnológico de Costa Rica\par}
\vspace{0.3cm}
{\large Escuela de Ingeniería Electrónica\par}
\vspace{2.5cm}

{\color{azulunc}\rule{\textwidth}{1.5pt}}
\vspace{0.6cm}

{\Huge\bfseries Acelerador de Conversión\\[0.2cm] RGB a Escala de Grises\par}
\vspace{0.5cm}
{\Large Implementación en HLS y Prototipo Virtual sobre gem5\par}

\vspace{0.6cm}
{\color{azulunc}\rule{\textwidth}{1.5pt}}

\vspace{2cm}

{\large\textbf{Evaluación Corta 2}\par}
\vspace{0.3cm}
{\large Diseño de Alto Nivel --- MP6160\par}
{\large II Cuatrimestre 2026\par}

\vspace{2cm}

\begin{tabular}{r l}
\textbf{Profesor:} & Luis G. León-Vega, Ph.D. \\[0.35cm]
\textbf{Estudiantes:} & \rule{6cm}{0.4pt} \\[0.35cm]
                      & \rule{6cm}{0.4pt} \\[0.35cm]
                      & \rule{6cm}{0.4pt} \\
\end{tabular}

\vfill
{\large \today\par}
\end{titlepage}

%-------------------------------- INDICE --------------------------------------
\tableofcontents
\newpage

%==============================================================================
\section{Introducción}
%==============================================================================

El presente reporte documenta el desarrollo de un acelerador de procesamiento de
imagen que convierte una imagen en formato RAW RGB de resolución 1080p a escala
de grises. El trabajo parte del modelo a nivel de transacciones desarrollado en
la evaluación anterior (SystemC) y lo proyecta en dos direcciones
complementarias:

\begin{enumerate}[leftmargin=*]
  \item Una \textbf{implementación de hardware mediante síntesis de alto nivel}
        (HLS) con Vitis 2024.1, con la tarjeta AMD Kria KV260 como objetivo.
  \item Un \textbf{prototipo virtual} basado en ARM64, construido sobre el
        simulador gem5, donde el acelerador se integra como periférico
        \emph{memory-mapped} mediante TLM-2.0 y es controlado por un programa
        escrito en C.
\end{enumerate}

Ambas implementaciones derivan del mismo modelo de referencia, por lo que deben
producir resultados funcionalmente equivalentes. La Parte~\ref{part:correspondencia}
documenta la verificación cruzada que lo comprueba.

\subsection{Flujo funcional del sistema}

El sistema implementa el siguiente flujo, heredado de la evaluación anterior:

\begin{enumerate}[leftmargin=*,itemsep=2pt]
  \item Cargar una imagen desde un almacenamiento persistente.
  \item La imagen está en formato RAW RGB, resolución 1080p
        ($1920 \times 1080$).
  \item Indicar al acelerador la dirección base de entrada, la de salida y la
        cantidad de píxeles a procesar.
  \item El acelerador convierte la imagen a escala de grises.
  \item El resultado se almacena nuevamente en disco.
\end{enumerate}

\subsection{Conversión implementada}

La conversión emplea la luminancia BT.601 evaluada en \textbf{aritmética
entera}, idéntica en ambas implementaciones:

\begin{equation}
Y = \frac{77 \cdot R + 150 \cdot G + 29 \cdot B}{256}
  \equiv (77R + 150G + 29B) \gg 8
\label{eq:bt601}
\end{equation}

Los coeficientes $77$, $150$ y $29$ se obtienen al redondear
$0{,}299 \times 256$, $0{,}587 \times 256$ y $0{,}114 \times 256$
respectivamente. El uso de aritmética entera
--- y no de punto flotante --- es un requisito para que ambas implementaciones
sean equivalentes bit a bit, según se detalla en la
Sección~\ref{sec:bitexact}.

\subsection{Parámetros de la imagen}

\begin{table}[H]
\centering
\caption{Dimensiones de los datos procesados}
\begin{tabular}{@{}lr@{}}
\toprule
\textbf{Parámetro} & \textbf{Valor} \\
\midrule
Resolución                      & $1920 \times 1080$ \\
Total de píxeles                & 2\,073\,600 \\
Bytes por píxel (entrada RGB)   & 3 \\
Bytes por píxel (salida gris)   & 1 \\
Tamaño de la imagen de entrada  & 6\,220\,800 bytes \\
Tamaño de la imagen de salida   & 2\,073\,600 bytes \\
\bottomrule
\end{tabular}
\end{table}

\newpage
%==============================================================================
\part{Implementación en HLS}
\label{part:hls}
%==============================================================================

\pendiente{Esta parte completa del documento debe ser desarrollada por el equipo
encargado de la implementación en HLS. A continuación se proponen las secciones
que cubren los requisitos del enunciado. Cada una puede ampliarse, reordenarse o
renombrarse según convenga.}

%------------------------------------------------------------------------------
\section{Arquitectura del acelerador en HLS}
%------------------------------------------------------------------------------

\pendiente{Describir la arquitectura propuesta: interfaces empleadas
(AXI4 Memory-Mapped o AXI4-Stream para datos, AXI4-Lite para control),
organización de los módulos y decisiones de diseño.}

\subsection{Interfaces de comunicación}

\pendiente{Documentar las interfaces AXI4 y AXI4-Lite: puertos, anchos de bus y
\emph{pragmas} empleados.}

\subsection{Estructura del \emph{pipeline}}

\pendiente{Mostrar la separación explícita entre las etapas de entrada/salida de
datos y la etapa de procesamiento, tal como exige el enunciado. Incluir el
fragmento de código donde se aprecia esta separación y los \emph{pragmas} de
\texttt{PIPELINE} / \texttt{DATAFLOW} utilizados.}

\subsection{Diagrama de bloques de la implementación}

\pendiente{Insertar el diagrama de bloques de la arquitectura en hardware.}

% Ejemplo de como insertar una figura:
%
% \begin{figure}[H]
%   \centering
%   \includegraphics[width=0.85\textwidth]{figuras/diagrama_hls.png}
%   \caption{Arquitectura del acelerador implementado en HLS.}
%   \label{fig:diagrama-hls}
% \end{figure}

%------------------------------------------------------------------------------
\section{Testbench y co-simulación}
%------------------------------------------------------------------------------

\pendiente{Describir el testbench en C empleado para la co-simulación C/RTL:
cómo se genera o carga la imagen de prueba, cómo se compara el resultado y qué
criterio de aceptación se utiliza.}

% \begin{lstlisting}[style=cstyle,caption={Testbench en C para co-simulacion.}]
% // ... codigo del testbench ...
% \end{lstlisting}

%------------------------------------------------------------------------------
\section{Flujo de síntesis y scripts TCL}
%------------------------------------------------------------------------------

\pendiente{Documentar los scripts TCL que automatizan el flujo completo de HLS
(creación del proyecto, C-Sim, síntesis, co-simulación y exportación del IP), e
indicar cómo ejecutarlos.}

% \begin{lstlisting}[style=shstyle,caption={Ejecucion del flujo de HLS.}]
% vitis_hls -f scripts/run_hls.tcl
% \end{lstlisting}

%------------------------------------------------------------------------------
\section{Resultados de síntesis}
\label{sec:resultados-hls}
%------------------------------------------------------------------------------

\pendiente{Completar con los datos del reporte de síntesis. Estos valores son
necesarios además para calibrar el modelo de latencia del prototipo virtual
(véase la Sección~\ref{sec:latencia-cruzada}).}

\begin{table}[H]
\centering
\caption{Resultados de síntesis en Vitis HLS 2024.1 (Kria KV260)}
\label{tab:sintesis}
\begin{tabular}{@{}ll@{}}
\toprule
\textbf{Métrica} & \textbf{Valor obtenido} \\
\midrule
Frecuencia objetivo          & 250 MHz \\
Frecuencia alcanzada         & \rule{4cm}{0.4pt} \\
Periodo estimado             & \rule{4cm}{0.4pt} \\
\emph{Initiation Interval} (II) & \rule{4cm}{0.4pt} \\
Latencia total (ciclos)      & \rule{4cm}{0.4pt} \\
Latencia total ($\mu$s)      & \rule{4cm}{0.4pt} \\
\midrule
LUT                          & \rule{4cm}{0.4pt} \\
FF                           & \rule{4cm}{0.4pt} \\
BRAM                         & \rule{4cm}{0.4pt} \\
DSP                          & \rule{4cm}{0.4pt} \\
\bottomrule
\end{tabular}
\end{table}

\subsection{Análisis de los resultados de síntesis}

\pendiente{Analizar el uso de recursos, el cumplimiento de la restricción
temporal de 250 MHz y el comportamiento del \emph{pipeline}.}

%------------------------------------------------------------------------------
\section{Mapa de registros AXI4-Lite}
\label{sec:mapa-axi}
%------------------------------------------------------------------------------

\pendiente{Documentar el mapa de registros generado por Vitis (disponible en el
archivo de cabecera del \emph{driver} exportado). Esta información se utiliza en
la Tabla~\ref{tab:correspondencia-registros} para establecer la correspondencia
con el prototipo virtual.}

\newpage
%==============================================================================
\part{Prototipo Virtual}
\label{part:vp}
%==============================================================================

Esta parte documenta la construcción, ejecución y verificación del prototipo
virtual: un sistema ARM64 simulado en gem5 sobre el cual se ejecuta un programa
en C que controla el acelerador SystemC a través de transacciones TLM-2.0.

%------------------------------------------------------------------------------
\section{Fundamento: qué es un prototipo virtual}
%------------------------------------------------------------------------------

Un \textbf{prototipo virtual} es un modelo ejecutable de un sistema de hardware,
suficientemente fiel como para que el software real se ejecute sobre él antes de
que el hardware exista físicamente. Su valor consiste en romper la dependencia
secuencial entre el desarrollo de hardware y el de software: el equipo de
software recibe un modelo ejecutable del sistema y comienza a trabajar de
inmediato.

El criterio que determina si un modelo merece esa denominación es directo:
\emph{¿funcionaría, sin modificaciones, el software escrito para el prototipo
sobre el chip real?} En esta implementación la respuesta es afirmativa, ya que
el programa en C no contiene ninguna construcción dependiente de la simulación:
únicamente lee y escribe direcciones de memoria, exactamente como lo haría un
controlador sobre silicio.

\subsection{gem5 como plataforma de simulación}

gem5 es un simulador de sistemas de computador: no simula un programa aislado,
sino una computadora completa --- procesador, jerarquía de memoria, buses y
periféricos --- con fidelidad suficiente para ejecutar un binario real. En este
trabajo aporta cuatro elementos:

\begin{itemize}[leftmargin=*,itemsep=2pt]
  \item Una CPU ARM64 (ARMv8-A) que ejecuta instrucciones ARM reales.
  \item Una plataforma completa (\texttt{VExpress\_GEM5}) con buses,
        controlador de interrupciones y UART.
  \item La capacidad de conectar componentes propios al bus del sistema, que es
        el mecanismo empleado para integrar el acelerador.
  \item Un núcleo de SystemC integrado (\texttt{SystemC\_Kernel}) que permite
        que los módulos SystemC y los componentes de gem5 avancen con el mismo
        planificador de eventos.
\end{itemize}

\subsection{Justificación del nivel de abstracción TLM-2.0}

El acelerador podría modelarse a nivel de señales (RTL), pero ello resultaría
órdenes de magnitud más lento sin aportar información adicional para la
validación del software: al programa en C no le concierne cómo evolucionan las
señales internas, sino si al escribir en el registro de control el acelerador
inicia su operación. TLM-2.0 modela la comunicación como \emph{transacciones}
--- objetos que encapsulan comando, dirección, datos y longitud --- lo que
constituye el nivel de abstracción adecuado para un prototipo virtual: rápido y
exacto en aquello que resulta relevante.

%------------------------------------------------------------------------------
\section{Arquitectura del prototipo virtual}
%------------------------------------------------------------------------------

La totalidad del sistema se ejecuta dentro de un único binario de gem5, bajo el
planificador de \texttt{SystemC\_Kernel}. El camino que recorre un acceso desde
la CPU hasta el acelerador es el siguiente:

\begin{lstlisting}[style=diagrama,caption={Arquitectura del prototipo virtual.}]
   +===========================================================================+
   |                        gem5  (SystemC_Kernel)                             |
   |   +---------------------+                                                 |
   |   |   CPU ARM64 (ARMv8) |   ejecuta sw/accel_app.elf (bare-metal, C)      |
   |   +----------+----------+                                                 |
   |              | acceso MMIO                                                |
   |        +-----v---------------------------------------------------+        |
   |        |                 membus  (SystemXBar)                    |        |
   |        +--+-------------------+-------------------------+--------+        |
   |           v                   v                         v                 |
   |   +-------+------+    +-------+--------+     +----------+-----------+      |
   |   | DRAM         |    | UART PL011     |     | iobridge --> iobus   |      |
   |   | 0x80000000   |    | 0x1c090000     |     |          |           |      |
   |   +--------------+    +----------------+     +----------v-----------+      |
   |                                              | Gem5ToTlmBridge32    |      |
   |                                              | rango: 0x2F000000    |      |
   |                                              +----------+-----------+      |
   |                                                         | TLM-2.0          |
   |. . . . . . . . . . frontera gem5 <-> SystemC . . . . . .v . . . . . . . . .|
   |                                              +----------+-----------+      |
   |                                              |   ImgAccel (SystemC) |      |
   |                                              |  tSocket + wrapper   |      |
   |                                              |  b_transport()       |      |
   |                                              |  SC_THREAD proceso() |      |
   |                                              +----+-------------+---+      |
   +===================================================|=============|==========+
                                              lee      v             v  escribe
                                          +------------+--+   +------+---------+
                                          | input .rgb    |   | output .raw    |
                                          +---------------+   +----------------+
                                            Almacenamiento persistente (host)
\end{lstlisting}

Una característica relevante del diseño es que el acelerador constituye un
\textbf{\emph{target} puro}: el tráfico TLM circula siempre en sentido CPU
$\rightarrow$ acelerador. El movimiento de los píxeles lo efectúa el propio
acelerador contra archivos del sistema anfitrión, de modo que no se requiere un
\emph{transactor} maestro de retorno hacia gem5.

\subsection{Organización del repositorio}

\begin{table}[H]
\centering
\caption{Organización de los archivos del prototipo virtual}
\begin{tabular}{@{}p{4.9cm}p{9.3cm}@{}}
\toprule
\textbf{Ruta} & \textbf{Contenido} \\
\midrule
\multicolumn{2}{@{}l}{\itshape Implementación en HLS} \\
\quad\texttt{hls/} & Fuentes sintetizables, testbench, scripts TCL y reportes
  de síntesis. \\[2pt]
\multicolumn{2}{@{}l}{\itshape Prototipo virtual --- \texttt{virtual\_prototype/}} \\
\quad\texttt{systemc/} & Modelo del acelerador: contrato, modelo SystemC y
  envoltorio para gem5. \\
\quad\texttt{sw/} & Programa en C \emph{bare-metal}, arranque, \emph{linker
  script} y \emph{Makefile}. \\
\quad\texttt{gem5-config/} & Configuración del SoC: construye el sistema ARM64
  y conecta el acelerador. \\
\quad\texttt{verify/} & Banco de pruebas independiente en SystemC puro. \\
\quad\texttt{scripts/} & Automatización: verificación, compilación y
  ejecución. \\[2pt]
\multicolumn{2}{@{}l}{\itshape Compartido} \\
\quad\texttt{images/} & Imagen RAW RGB de entrada e imagen de salida. \\
\quad\texttt{docs/} & Diagramas, salida de referencia y este reporte. \\
\quad\texttt{scripts/} & Verificación cruzada entre ambas implementaciones. \\
\bottomrule
\end{tabular}
\end{table}

\subsection{Módulos del acelerador}

El directorio \texttt{virtual\_prototype/systemc/img\_accel\_vp/} contiene seis
organizados en tres capas funcionales:

\begin{table}[H]
\centering
\caption{Módulos que componen el acelerador}
\begin{tabular}{@{}p{3.5cm}p{2.2cm}p{8.5cm}@{}}
\toprule
\textbf{Archivo} & \textbf{Capa} & \textbf{Función} \\
\midrule
\texttt{accel\_core.h}    & Contrato   & Constantes puras: mapa de registros,
  dimensiones y la función de conversión. Sin dependencias de SystemC. \\
\texttt{accel\_io.h}      & Contrato   & Lectura y escritura de archivos:
  modela el almacenamiento persistente. \\
\texttt{sc\_img\_accel.hh}& Modelo     & Declaración del \texttt{SC\_MODULE}:
  socket TLM, registros internos y evento de arranque. \\
\texttt{sc\_img\_accel.cc}& Modelo     & Implementación de \texttt{b\_transport()}
  y del \texttt{SC\_THREAD} de conversión. \\
\texttt{ImgAccel.py}      & Envoltorio & Descripción del acelerador como
  \emph{SimObject} de gem5. \\
\texttt{SConscript}       & Envoltorio & Regla de compilación para el mecanismo
  \texttt{EXTRAS} de gem5. \\
\bottomrule
\end{tabular}
\end{table}

La separación entre \emph{contrato} y \emph{modelo} obedece a un objetivo
concreto: el banco de pruebas independiente incluye el mismo
\texttt{accel\_core.h}, de manera que lo verificado fuera de gem5 es
literalmente el mismo código que se ejecuta dentro del prototipo.

%------------------------------------------------------------------------------
\section{Mapa de memoria y registros}
\label{sec:mapa-vp}
%------------------------------------------------------------------------------

\begin{table}[H]
\centering
\caption{Mapa de memoria del SoC simulado}
\begin{tabular}{@{}lll@{}}
\toprule
\textbf{Rango} & \textbf{Periférico} & \textbf{Función} \\
\midrule
\texttt{0x1C090000}                 & UART PL011 & Consola del programa en C \\
\texttt{0x2F000000 -- 0x2F000FFF}   & \texttt{ImgAccel} & Ventana de registros
  del acelerador \\
\texttt{0x80000000 \ldots}          & DRAM & Memoria principal; carga y ejecuta
  el ELF \\
\bottomrule
\end{tabular}
\end{table}

\begin{table}[H]
\centering
\caption{Registros del acelerador (base \texttt{0x2F000000})}
\label{tab:registros-vp}
\begin{tabular}{@{}llll@{}}
\toprule
\textbf{Offset} & \textbf{Nombre} & \textbf{Acceso} & \textbf{Descripción} \\
\midrule
\texttt{0x00} & \texttt{DIR\_IN}     & RW & Base lógica de la imagen de entrada \\
\texttt{0x04} & \texttt{DIR\_OUT}    & RW & Base lógica de la imagen de salida \\
\texttt{0x08} & \texttt{NUM\_PIXELS} & RW & Píxeles a procesar (2\,073\,600) \\
\texttt{0x0C} & \texttt{CTRL}        & WO & bit 0 = \texttt{START} \\
\texttt{0x10} & \texttt{STATUS}      & RO & bit 0 = \texttt{BUSY},
                                            bit 1 = \texttt{DONE} \\
\texttt{0x14} & \texttt{ID}          & RO & Identificador \texttt{0xACCE1111} \\
\bottomrule
\end{tabular}
\end{table}

La dirección base se encuentra definida en dos ubicaciones que deben mantenerse
consistentes: \texttt{virtual\_prototype/sw/accel.h} (software) y
\texttt{virtual\_prototype/gem5-config/vp\_arm.py} (sistema). El registro de
identificación permite al programa detectar una eventual inconsistencia entre
ambas.

%------------------------------------------------------------------------------
\section{Formato de las transacciones}
%------------------------------------------------------------------------------

Cada acceso MMIO de la CPU a la ventana del acelerador genera una transacción
\texttt{tlm\_generic\_payload} de 4~bytes. El puente
\texttt{Gem5ToTlmBridge32} la entrega al método \texttt{b\_transport()}
del módulo SystemC.

\begin{table}[H]
\centering
\caption{Campos del \emph{payload} TLM-2.0}
\begin{tabular}{@{}p{3.3cm}p{5.3cm}p{6.2cm}@{}}
\toprule
\textbf{Campo} & \textbf{Valor} & \textbf{Descripción} \\
\midrule
\texttt{command}          & \texttt{TLM\_READ/WRITE\_COMMAND} & Según sea
  lectura o escritura \\
\texttt{address}          & Offset dentro de \texttt{0x2F000000} & El puente
  entrega el offset local \\
\texttt{data\_ptr}        & Puntero a 4 bytes & Valor del registro \\
\texttt{data\_length}     & \texttt{4} & Registro de 32 bits \\
\texttt{streaming\_width} & \texttt{4} & Sin \emph{streaming} \\
\texttt{byte\_enable\_ptr}& \texttt{nullptr} & No se emplean \emph{byte enables} \\
\texttt{response\_status} & \texttt{TLM\_OK\_RESPONSE} & Error si
  \texttt{data\_length} $\neq$ 4 \\
\bottomrule
\end{tabular}
\end{table}

\subsection{Secuencia de transacciones}

\begin{table}[H]
\centering
\caption{Transacciones generadas por el programa en C}
\begin{tabular}{@{}lllll@{}}
\toprule
\textbf{Paso} & \textbf{Comando} & \textbf{Registro} & \textbf{Valor} &
\textbf{Sentido} \\
\midrule
Identificar & READ  & \texttt{ID} (\texttt{0x14})         & \texttt{0xACCE1111}
  & acel. $\rightarrow$ CPU \\
Programar   & WRITE & \texttt{DIR\_IN} (\texttt{0x00})    & \texttt{0}
  & CPU $\rightarrow$ acel. \\
Programar   & WRITE & \texttt{DIR\_OUT} (\texttt{0x04})   & \texttt{0}
  & CPU $\rightarrow$ acel. \\
Programar   & WRITE & \texttt{NUM\_PIXELS} (\texttt{0x08})& 2\,073\,600
  & CPU $\rightarrow$ acel. \\
Arrancar    & WRITE & \texttt{CTRL} (\texttt{0x0C})       & \texttt{0x1}
  & CPU $\rightarrow$ acel. \\
Sondear     & READ  & \texttt{STATUS} (\texttt{0x10})     &
  \texttt{BUSY}/\texttt{DONE} & acel. $\rightarrow$ CPU \\
\bottomrule
\end{tabular}
\end{table}

%------------------------------------------------------------------------------
\section{Diagrama de secuencias}
%------------------------------------------------------------------------------

\begin{lstlisting}[style=diagrama,caption={Secuencia de operacion del prototipo virtual.}]
 CPU ARM64 (main.c)      membus        Gem5ToTlmBridge32       ImgAccel (SystemC)
     |                     |                  |                        |
     |  MMIO RD 0x2F...14  |                  |                        |
     |-------------------->|--- TLM READ ---->|----- b_transport ----->|
     |                     |                  |     ID = 0xACCE1111    |
     |<--------------------|<-----------------|<-----------------------|
     |  (verifica ID)      |                  |                        |
     |                     |                  |                        |
     |  MMIO WR DIR_IN / DIR_OUT / NUM_PIXELS |                        |
     |-------------------->|--- TLM WRITE --->|----- b_transport ----->| (guarda)
     |                     |                  |                        |
     |  MMIO WR CTRL=START |                  |                        |
     |-------------------->|--- TLM WRITE --->|----- b_transport ----->| notify()
     |<-- OK --------------|<-----------------|<-----------------------|    |
     |                     |                  |          +-------------v-----------+
     |  bucle:             |                  |          | SC_THREAD proceso():    |
     |   MMIO RD STATUS    |                  |          |  cargar input .rgb      |
     |-------------------->|--- TLM READ ---->|--------->|  rgb_a_gris (BT.601)    |
     |<-- BUSY/DONE -------|<-----------------|<---------|  wait(npix*4 ns)        |
     |   (repite hasta DONE)                  |          |  guardar output .raw    |
     |                     |                  |          |  done = true            |
     |                     |                  |          +-------------------------+
     |  STATUS & DONE == 1 |                  |                        |
     |  -> EOT al UART -> gem5 finaliza       |                        |
\end{lstlisting}

%------------------------------------------------------------------------------
\section{El programa en C}
%------------------------------------------------------------------------------

El software se ejecuta en modo \emph{bare-metal}: sin sistema operativo, sin
biblioteca estándar y con la MMU desactivada, de modo que las direcciones son
físicas. Esta condición explica tres características del código: la existencia
de un controlador propio del UART, de una rutina de arranque en ensamblador y de
un \emph{linker script} explícito.

\subsection{Acceso a los registros}

El acceso al periférico se realiza mediante punteros calificados con
\texttt{volatile}:

\begin{lstlisting}[style=cstyle,caption={Accesores MMIO (\texttt{sw/accel.h}).}]
#define ACCEL_BASE   0x2f000000UL

static inline void accel_wr(uint32_t off, uint32_t val)
{
    *(volatile uint32_t *)(ACCEL_BASE + off) = val;
}

static inline uint32_t accel_rd(uint32_t off)
{
    return *(volatile uint32_t *)(ACCEL_BASE + off);
}
\end{lstlisting}

El calificador \texttt{volatile} resulta indispensable. El compilador, al
optimizar, presupone que el contenido de la memoria únicamente se modifica
cuando el propio programa la escribe; bajo esa premisa extraería la lectura del
registro de estado fuera del bucle de espera, convirtiéndolo en un bucle
infinito. Al declarar el acceso como \texttt{volatile}, el compilador queda
obligado a emitir una lectura efectiva del bus en cada iteración.

\subsection{Flujo del programa principal}

\begin{lstlisting}[style=cstyle,caption={Programa principal (\texttt{sw/main.c}), estructura.}]
int main(void)
{
    /* 1) Identificar el periferico */
    uint32_t id = accel_rd(REG_ID);
    if (id != MAGIC_ID) {
        uart_puts("[SW] ERROR: periferico no encontrado\n");
        return 1;
    }

    /* 2) Programar los registros de configuracion */
    accel_wr(REG_DIR_IN,     0);
    accel_wr(REG_DIR_OUT,    0);
    accel_wr(REG_NUM_PIXELS, IMG_PIXELS);   /* 1920 x 1080 */

    /* 3) Arrancar el acelerador */
    accel_wr(REG_CTRL, CTRL_START);

    /* 4) Sondear STATUS hasta que termine */
    uint32_t status;
    do {
        status = accel_rd(REG_STATUS);
    } while (!(status & ST_DONE));

    uart_puts("[SW] DONE.\n");
    return 0;
}
\end{lstlisting}

La escritura del bit \texttt{START} retorna de forma inmediata: el método
\texttt{b\_transport()} únicamente notifica un evento, mientras que el
procesamiento se realiza en un \texttt{SC\_THREAD} independiente. Esto permite
que la CPU continúe su ejecución y sondee el estado, reproduciendo el
comportamiento de un periférico real que opera en paralelo con el procesador.

\subsection{Del código C a las transacciones TLM}

El desensamblado del binario permite observar la correspondencia directa entre
el código fuente y el tráfico generado:

\begin{lstlisting}[style=consola,caption={Bucle de sondeo: codigo C y ensamblador ARM64 generado.}]
/* C */   accel_wr(REG_CTRL, CTRL_START);
          do { status = accel_rd(REG_STATUS); } while (!(status & ST_DONE));

/* ASM generado */
  mov   w1, #0x1
  str   w1, [x6], #4      // escribe START en CTRL (0x..0c); x6 pasa a 0x..10
  ...
  ldr   w0, [x6]          // lee STATUS (0x2f000010)  <-- transaccion TLM
  tbz   w0, #1, <volver>  // si el bit DONE es 0, repite
\end{lstlisting}

Se aprecia que la instrucción \texttt{ldr} permanece \emph{dentro} del bucle: es
la evidencia, en el binario generado, de que el calificador \texttt{volatile}
cumplió su función.

%------------------------------------------------------------------------------
\section{Construcción y ejecución}
%------------------------------------------------------------------------------

\subsection{Requisitos}

\begin{lstlisting}[style=shstyle,caption={Dependencias del entorno.}]
# SystemC
sudo apt install -y libsystemc-dev g++
# Cross-compiler ARM64
sudo apt install -y gcc-aarch64-linux-gnu
# Dependencias de gem5
sudo apt install -y build-essential scons python3-dev git m4 \
    zlib1g-dev libprotobuf-dev protobuf-compiler libboost-all-dev pkg-config
\end{lstlisting}

\subsection{Procedimiento}

\begin{lstlisting}[style=shstyle,caption={Construccion y ejecucion del prototipo virtual.}]
cd virtual_prototype

# 1) Verificar el acelerador (sin gem5)
./scripts/verify.sh

# 2) Compilar el programa en C -> sw/accel_app.elf
./scripts/build_sw.sh

# 3) Compilar gem5 con el acelerador integrado (EXTRAS)
yes | ./scripts/build_gem5.sh

# 4) Ejecutar el prototipo virtual completo
./scripts/run_vp.sh
\end{lstlisting}

El paso 3 emplea el mecanismo \texttt{EXTRAS} de gem5, que permite compilar
código externo dentro del binario del simulador sin modificar su árbol de
fuentes:

\begin{lstlisting}[style=shstyle]
scons build/ARM/gem5.opt EXTRAS=<repo>/systemc/img_accel_vp -j$(nproc)
\end{lstlisting}

\newpage
%------------------------------------------------------------------------------
\section{Resultados}
\label{sec:resultados-vp}
%------------------------------------------------------------------------------

\subsection{Entorno de ejecución}

\begin{table}[H]
\centering
\caption{Entorno en el que se obtuvieron los resultados}
\begin{tabular}{@{}ll@{}}
\toprule
\textbf{Componente} & \textbf{Versión} \\
\midrule
Sistema operativo   & Debian 13 (Trixie) \\
gem5                & 25.1.0.1 \\
SystemC             & 3.0.2 (Accellera) \\
Cross-compiler      & \texttt{aarch64-linux-gnu-gcc} 14.2.0 \\
CPU simulada        & \texttt{AtomicSimpleCPU}, ARMv8-A, 1 GHz \\
Plataforma          & \texttt{VExpress\_GEM5\_V2} \\
\bottomrule
\end{tabular}
\end{table}

\subsection{Verificación del acelerador de forma independiente}

Previo a la integración, el acelerador se verificó mediante un banco de pruebas
en SystemC puro que reproduce el mismo protocolo de registros que ejercerá la
CPU ARM64:

\begin{lstlisting}[style=consola,caption={Salida de \texttt{./scripts/verify.sh}.}]
        SystemC 3.0.2-Accellera
[CPU] ID leido = 0xacce1111
[CPU] registros programados (npix=2073600)
[CPU] START enviado, sondeando STATUS...
[CPU] DONE tras 9 sondeos (9 ms)
[TB] simulacion OK
>> Comparando con la referencia...
OK: la salida es identica bit a bit a la referencia.
\end{lstlisting}

El acelerador respondió con el identificador esperado, aceptó la configuración,
inició la operación al recibir \texttt{START} y señalizó su terminación tras
nueve sondeos del registro de estado.

\subsection{Compilación del programa en C}

\begin{table}[H]
\centering
\caption{Atributos del binario generado}
\begin{tabular}{@{}ll@{}}
\toprule
\textbf{Atributo} & \textbf{Valor} \\
\midrule
Formato / Arquitectura & ELF64 / AArch64 \\
Tipo                   & \texttt{EXEC} (ejecutable, \emph{bare-metal}) \\
Punto de entrada       & \texttt{0x80000000} \\
Intérprete dinámico    & Ninguno \\
\bottomrule
\end{tabular}
\end{table}

La ausencia de intérprete dinámico confirma que el binario es efectivamente
\emph{bare-metal} y puede ser cargado directamente por gem5.

\subsection{Ejecución del prototipo virtual completo}

\begin{lstlisting}[style=consola,caption={Salida de \texttt{./scripts/run\_vp.sh}.}]
gem5 Simulator System.  https://www.gem5.org
gem5 version 25.1.0.1
...
info: kernel located at: .../sw/accel_app.elf
[ImgAccel] periferico en linea (in=.../input_1080p.rgb,
                                out=.../output_1080p_gray.raw)
systemc_kernel.system.terminal: Listening for connections on port 3456
== Prototipo virtual ARM64 + acelerador SystemC: iniciando ==
== Simulacion terminada @ tick 1312000: UART received EOT ==
>> Salida generada: .../images/output/output_1080p_gray.raw
-rw-r--r-- 1 daniel daniel 2073600 .../output_1080p_gray.raw
\end{lstlisting}

La simulación finalizó de forma controlada mediante el byte \texttt{EOT}
(\texttt{0x04}) que la rutina de arranque envía al UART al retornar de
\texttt{main()}. Ello confirma que el programa completó su ejecución: alcanzó el
final del bucle de sondeo, lo que a su vez implica que el acelerador señalizó su
terminación.

\subsection{Comprobación de resultados}
\label{sec:comprobacion}

La verificación de correctitud se realizó comparando la imagen generada por el
prototipo virtual con la salida del modelo de referencia de la evaluación
anterior:

\begin{lstlisting}[style=consola,caption={Verificacion mediante suma de comprobacion.}]
$ md5sum images/output/output_1080p_gray.raw docs/reference_output_gray.raw

45041cfcf4f0f1f0d8cdc2230acee01b  images/output/output_1080p_gray.raw
45041cfcf4f0f1f0d8cdc2230acee01b  docs/reference_output_gray.raw
\end{lstlisting}

\begin{table}[H]
\centering
\caption{Resultados de la verificación}
\begin{tabular}{@{}lll@{}}
\toprule
\textbf{Criterio} & \textbf{Resultado} & \textbf{Estado} \\
\midrule
Tamaño de la salida   & 2\,073\,600 bytes & \textcolor{verdeok}{Correcto} \\
Suma MD5              & Coincide con la referencia &
  \textcolor{verdeok}{Correcto} \\
Bytes discrepantes    & 0 de 2\,073\,600 & \textcolor{verdeok}{Correcto} \\
Terminación           & \texttt{UART received EOT} &
  \textcolor{verdeok}{Correcto} \\
\bottomrule
\end{tabular}
\end{table}

La coincidencia de las sumas de comprobación indica que los 2\,073\,600 bytes de
ambas imágenes son idénticos. Se concluye que la migración del modelo a una
arquitectura ARM64 con comunicación TLM-2.0 preservó íntegramente la
funcionalidad del acelerador original.

\subsection{Análisis de los resultados}

\subsubsection{Correctitud funcional}

La equivalencia bit a bit constituye el criterio más estricto de verificación
disponible: no admite tolerancias ni aproximaciones. Su cumplimiento demuestra
que las adaptaciones introducidas al modelo --- el envoltorio del socket para el
puente de gem5 y la modificación del camino de datos --- no alteraron la
función de transferencia del acelerador.

\subsubsection{Modelo temporal}

El acelerador modela su latencia de procesamiento como
$t = N_{\text{px}} \times 4\,\text{ns}$, correspondiente a un píxel por ciclo a
250 MHz. Para la imagen completa:

\begin{equation}
t = 2\,073\,600 \times 4\,\text{ns} = 8{,}29\,\text{ms}
\end{equation}

Este valor asume un \emph{Initiation Interval} unitario en el
\emph{pipeline}. La Sección~\ref{sec:latencia-cruzada} contrasta esta
suposición con los resultados de síntesis.

\subsubsection{Validez del prototipo}

El programa en C no contiene ninguna construcción específica de la simulación:
únicamente accede a direcciones de memoria mediante punteros \texttt{volatile}.
En consecuencia, el mismo código fuente, recompilado, sería válido para
controlar el IP sintetizado sobre la tarjeta KV260 --- modificando únicamente la
dirección base del periférico. Esta propiedad es precisamente la que otorga
valor a un prototipo virtual.

%------------------------------------------------------------------------------
\section{Problemas de integración encontrados}
%------------------------------------------------------------------------------

Durante la puesta en marcha se presentaron dos incidencias cuya resolución
resulta ilustrativa del funcionamiento del entorno.

\subsection{Generación de un ejecutable con enlace dinámico}

Al ejecutar el prototipo por primera vez, gem5 reportó:

\begin{lstlisting}[style=consola]
fatal: Failed to open file /lib/ld-linux-aarch64.so.1.
\end{lstlisting}

\textbf{Causa.} El compilador GCC 14 genera de forma predeterminada ejecutables
independientes de posición (PIE), que requieren un intérprete dinámico
inexistente en un entorno \emph{bare-metal}. Cabe señalar que las opciones
\texttt{-nostdlib} y \texttt{-nostartfiles} no desactivan este comportamiento.

\textbf{Solución.} Incorporar las opciones \texttt{-fno-pie}, \texttt{-no-pie} y
\texttt{-static} al proceso de compilación. La verificación se realiza mediante
\texttt{readelf -l}, comprobando la ausencia de un segmento \texttt{INTERP}.

\subsection{Conflicto de rangos de direcciones}

\begin{lstlisting}[style=consola]
fatal: membus has two ports responding within range [0x2f000000:0x80000000]:
	system.iobridge.cpu_side_port
	system.tlm_bridge.gem5
\end{lstlisting}

\textbf{Causa.} El puente del acelerador se había conectado al \texttt{membus},
cuyo rango en dicha zona ya se encuentra cubierto por el \texttt{iobridge} de la
plataforma VExpress.

\textbf{Solución.} Conectar el puente al \texttt{iobus}, que es el bus destinado
a periféricos \emph{memory-mapped} en las plataformas ARM de gem5:

\begin{lstlisting}[style=pystyle]
system.tlm_bridge = Gem5ToTlmBridge32(gem5=system.iobus.mem_side_ports)
\end{lstlisting}

\newpage
%==============================================================================
\part{Correspondencia entre ambas implementaciones}
\label{part:correspondencia}
%==============================================================================

Ambas implementaciones derivan del mismo modelo de referencia, por lo que deben
resultar funcionalmente equivalentes. Esta parte documenta las correspondencias
y el procedimiento de verificación cruzada.

%------------------------------------------------------------------------------
\section{Equivalencia funcional}
\label{sec:bitexact}
%------------------------------------------------------------------------------

Las tres implementaciones deben producir la misma imagen de salida a partir de
la misma entrada:

\begin{lstlisting}[style=shstyle,caption={Verificacion cruzada de las tres salidas.}]
./scripts/cross_check.sh <salida_del_testbench_HLS.raw>
\end{lstlisting}

\begin{table}[H]
\centering
\caption{Verificación cruzada de equivalencia funcional}
\label{tab:bitexact}
\begin{tabular}{@{}lll@{}}
\toprule
\textbf{Implementación} & \textbf{Suma MD5} & \textbf{Estado} \\
\midrule
Modelo de referencia (Ev. 1) & \texttt{45041cfcf4f0f1f0d8cdc2230acee01b} &
  \textcolor{verdeok}{Verificado} \\
Prototipo virtual            & \texttt{45041cfcf4f0f1f0d8cdc2230acee01b} &
  \textcolor{verdeok}{Verificado} \\
Implementación HLS           & \rule{5cm}{0.4pt} & \emph{pendiente} \\
\bottomrule
\end{tabular}
\end{table}

\pendiente{Ejecutar \texttt{cross\_check.sh} con la salida del testbench de HLS
y completar la Tabla~\ref{tab:bitexact}.}

\subsection{Requisito de aritmética entera}

La equivalencia bit a bit únicamente se cumple si la implementación en HLS
emplea la misma aritmética entera de la Ecuación~\ref{eq:bt601}. El empleo de
tipos de punto flotante (\texttt{float}, \texttt{double}) o de coma fija
(\texttt{ap\_fixed}) introduce diferencias de redondeo.

\begin{table}[H]
\centering
\caption{Interpretación del diagnóstico de la verificación cruzada}
\begin{tabular}{@{}p{2.8cm}p{6.2cm}p{5cm}@{}}
\toprule
\textbf{Error abs. máx.} & \textbf{Diagnóstico} & \textbf{Acción} \\
\midrule
0     & Equivalencia bit a bit. & Ninguna. \\
1     & Diferencia de redondeo: se emplea punto flotante o coma fija. La lógica
        es correcta. & Migrar a aritmética entera. \\
2 -- 5& Coeficientes o política de redondeo distintos. & Revisar coeficientes y
        desplazamiento. \\
$>$ 10& Error lógico: canales intercambiados, fórmula distinta o
        desbordamiento. & Depurar el núcleo. \\
\bottomrule
\end{tabular}
\end{table}

Como referencia, una implementación en punto flotante con los coeficientes
decimales difiere del modelo entero en aproximadamente el 47\,\% de los píxeles,
pero siempre en $\pm 1$. Una proporción elevada de píxeles discrepantes no
implica necesariamente un error lógico.

%------------------------------------------------------------------------------
\section{Correspondencia de los mapas de registros}
%------------------------------------------------------------------------------

Ambas implementaciones exponen el mismo dispositivo desde la perspectiva del
software; únicamente difiere el mecanismo de transporte.

\begin{table}[H]
\centering
\caption{Correspondencia entre el prototipo virtual y la implementación en HLS}
\label{tab:correspondencia-registros}
\begin{tabular}{@{}p{4.2cm}p{4.6cm}p{5.2cm}@{}}
\toprule
\textbf{Función} & \textbf{Prototipo Virtual (TLM-2.0)} &
\textbf{HLS (AXI4-Lite)} \\
\midrule
Iniciar la operación      & \texttt{CTRL} (\texttt{0x0C}), bit 0 &
  \texttt{ap\_start} \\
Consultar terminación     & \texttt{STATUS} (\texttt{0x10}), bit 1 &
  \texttt{ap\_done} / \texttt{ap\_idle} \\
Indicar ocupado           & \texttt{STATUS} (\texttt{0x10}), bit 0 &
  \texttt{ap\_ready} / \texttt{ap\_idle} \\
Dirección base de entrada & \texttt{DIR\_IN} (\texttt{0x00}) &
  Puerto AXI4 Master \\
Dirección base de salida  & \texttt{DIR\_OUT} (\texttt{0x04}) &
  Puerto AXI4 Master \\
Cantidad de píxeles       & \texttt{NUM\_PIXELS} (\texttt{0x08}) &
  Argumento escalar AXI4-Lite \\
Identificación            & \texttt{ID} (\texttt{0x14}) & No aplica \\
\bottomrule
\end{tabular}
\end{table}

El protocolo de control es equivalente en ambos casos: escribir un bit de
arranque y sondear un bit de terminación. En consecuencia, el programa en C
desarrollado para el prototipo virtual constituye, conceptualmente, el mismo
controlador que se requeriría para operar el IP sobre la tarjeta KV260.

%------------------------------------------------------------------------------
\section{Calibración del modelo de latencia}
\label{sec:latencia-cruzada}
%------------------------------------------------------------------------------

El prototipo virtual modela la latencia del acelerador como:

\begin{lstlisting}[style=cppstyle]
wait(npix * 4.0, sc_core::SC_NS);   // 1 pixel/ciclo a 250 MHz
\end{lstlisting}

Dicho valor asume un \emph{Initiation Interval} unitario. Los resultados de
síntesis (Tabla~\ref{tab:sintesis}) permiten validar o corregir esta suposición:

\begin{table}[H]
\centering
\caption{Contraste entre el modelo temporal y la síntesis}
\begin{tabular}{@{}lll@{}}
\toprule
\textbf{Parámetro} & \textbf{Modelo del VP} & \textbf{Síntesis HLS} \\
\midrule
Frecuencia            & 250 MHz  & \rule{3.5cm}{0.4pt} \\
\emph{Init. Interval} & 1        & \rule{3.5cm}{0.4pt} \\
Latencia estimada     & 8,29 ms  & \rule{3.5cm}{0.4pt} \\
\bottomrule
\end{tabular}
\end{table}

\pendiente{Completar con los datos del reporte de síntesis. En caso de que el
\emph{Initiation Interval} difiera de la unidad, el modelo del prototipo virtual
debe ajustarse según $t = N_{\text{px}} \cdot II \cdot T_{\text{clk}}$.}

%------------------------------------------------------------------------------
\section{Diferencia arquitectónica en el camino de datos}
%------------------------------------------------------------------------------

Ambas implementaciones difieren en el tratamiento del movimiento de los píxeles,
diferencia que se documenta explícitamente:

\begin{table}[H]
\centering
\caption{Comparación del camino de datos}
\begin{tabular}{@{}p{2.6cm}p{5.6cm}p{5.8cm}@{}}
\toprule
 & \textbf{Implementación HLS} & \textbf{Prototipo Virtual} \\
\midrule
Control & AXI4-Lite
        & MMIO $\rightarrow$ TLM-2.0 \\
Datos   & AXI4 Master: el IP accede a la memoria del sistema
        & E/S directa del acelerador contra archivos del anfitrión \\
Sockets & Dos: control (esclavo) y datos (maestro)
        & Uno: el acelerador es \emph{target} puro \\
\bottomrule
\end{tabular}
\end{table}

\textbf{Justificación.} En el prototipo virtual el sistema ARM64 se ejecuta en
modo \emph{bare-metal}, sin sistema de archivos, por lo que la CPU no puede
cargar la imagen desde disco hacia la memoria sin habilitar \emph{semihosting}.
Se optó por modelar el almacenamiento persistente como entrada/salida del propio
acelerador --- de forma coherente con el rol que ya desempeñaba el módulo de
almacenamiento en la evaluación anterior --- lo que concentra el tráfico TLM en
el control del dispositivo. El enunciado autoriza expresamente adaptaciones del
modelo, y el resultado es funcionalmente idéntico según demuestra la
verificación de la Sección~\ref{sec:comprobacion}.

\newpage
%==============================================================================
\part*{Conclusiones}
\addcontentsline{toc}{part}{Conclusiones}
%==============================================================================
\section*{Conclusiones}
\addcontentsline{toc}{section}{Conclusiones}

\begin{enumerate}[leftmargin=*,itemsep=4pt]
  \item Se construyó un prototipo virtual funcional basado en ARM64 sobre gem5,
        en el cual un programa escrito en C, ejecutado en modo \emph{bare-metal}
        sobre una CPU ARMv8-A simulada, controla un acelerador modelado en
        SystemC mediante transacciones TLM-2.0.

  \item La verificación mediante sumas de comprobación demostró que la imagen
        producida por el prototipo virtual es idéntica, byte a byte, a la del
        modelo de referencia de la evaluación anterior, lo que confirma que la
        migración preservó la funcionalidad del acelerador.

  \item La integración del modelo SystemC dentro de gem5 mediante el mecanismo
        \texttt{EXTRAS} y el puente \texttt{Gem5ToTlmBridge32} permitió mantener
        el código del acelerador en el repositorio del proyecto, sin necesidad
        de modificar el árbol de fuentes del simulador.

  \item El software desarrollado no contiene construcciones dependientes de la
        simulación, por lo que resultaría válido para controlar el IP
        sintetizado sobre hardware real, modificando únicamente la dirección
        base del periférico. Esta propiedad constituye el criterio que valida el
        prototipo como tal.

  \pendiente{Añadir las conclusiones correspondientes a la implementación en HLS
  y a la comparación entre ambas.}
\end{enumerate}

\newpage
%==============================================================================
\section*{Declaración sobre el uso de Inteligencia Artificial}
\addcontentsline{toc}{section}{Declaración sobre el uso de IA}
%==============================================================================

De acuerdo con lo establecido en el enunciado, se declara el uso de herramientas
de Inteligencia Artificial durante el desarrollo de esta evaluación.

\subsection*{Parte de prototipo virtual}

Se utilizó \textbf{Claude (Anthropic)} a través de su interfaz de conversación,
con los siguientes propósitos:

\begin{itemize}[leftmargin=*,itemsep=3pt]
  \item \textbf{Consulta de conceptos y de API de integración.} Se consultó el
        procedimiento para integrar un módulo SystemC como periférico TLM-2.0
        dentro de un SoC ARM64 en gem5.
        \emph{Prompt representativo:} «cómo conectar un \emph{target} SystemC
        TLM-2.0 al bus de un sistema ARM64 en gem5 mediante
        \texttt{Gem5ToTlmBridge} y \texttt{SystemC\_Kernel}».

  \item \textbf{Generación de código base.} A partir del acelerador de la
        evaluación anterior, se generó con asistencia el andamiaje del módulo
        integrado a gem5 (\texttt{sc\_img\_accel.hh/.cc}, \texttt{ImgAccel.py},
        \texttt{SConscript}), el programa en C \emph{bare-metal}
        (\texttt{main.c}, \texttt{startup.S}, \texttt{link.ld}), la
        configuración de gem5 y los scripts de automatización.
        \emph{Prompt representativo:} «adaptar el acelerador RGB a grises de
        SystemC para que una CPU ARM64 lo controle por MMIO, y escribir el
        programa en C que lo maneja».

  \item \textbf{Depuración.} Se utilizó asistencia para diagnosticar los dos
        problemas de integración documentados: la generación de un ejecutable
        con enlace dinámico y el conflicto de rangos de direcciones entre el
        \texttt{membus} y el \texttt{iobridge}.

  \item \textbf{Generación de diagramas y mejora de la redacción.} Los diagramas
        de bloques y de secuencias en formato de texto, así como la redacción de
        la documentación técnica, se elaboraron con asistencia a partir del
        código del proyecto.
\end{itemize}

La totalidad del código fue revisada, integrada y verificada por el equipo. La
fórmula de conversión y el mapa de registros provienen del trabajo propio
desarrollado en la evaluación anterior.

\subsection*{Parte de implementación en HLS}

\pendiente{El equipo encargado de la implementación en HLS debe completar esta
subsección indicando si empleó herramientas de Inteligencia Artificial, con qué
propósito y los \emph{prompts} representativos. El enunciado establece que la
omisión de esta declaración implica la aplicación de la normativa de plagio; si
no se utilizaron dichas herramientas, debe indicarse explícitamente.}

\end{document}
